The previous report of this project (Modeling Dynamics) compared eight forecasting models on the six links and selected the Seq2Seq-GRU with additive attention. This report changes the data pipeline, the output of the model, the way forecasts are evaluated, and the way they are used. Table~\ref{tab:improved} lists every change against the previous report together with its measured effect. All numbers are on the test sessions, which no model saw during training, model selection, or calibration.

\begin{table}[tp]
\centering
\caption{What changed since the previous report and what each change achieved.}
\label{tab:improved}
\small
\begin{tabularx}{\linewidth}{>{\raggedright\arraybackslash}p{2.5cm}>{\raggedright\arraybackslash}X>{\raggedright\arraybackslash}X>{\raggedright\arraybackslash}X}
\toprule
Aspect & Previous report & This report & Effect \\
\midrule
Causality of the inputs & Missing Starlink pings were interpolated toward the next real ping (\InterpFracSLDL\% of the DL test targets), which leaks part of the future into the input. & Inputs hold the last real ping. Targets stay the measured series. & The leak hid \AuditPersHonePct\% of the 0.2\,s persistence RMSE. A model trained and tested on leaky data has a 0.2\,s error \LeakOptDSMAEhOne\% higher in deployment than reported (Section~\ref{sec:leakmodel}). \\
Horizon & 15 steps labeled as 15\,s. & 15 steps of 0.2\,s, i.e.\ 3\,s. & Correct reading of every horizon. \\
Uncertainty & 95\% intervals from the quantiles of the test residuals, so the coverage is 0.95 by construction. & Intervals calibrated on held-out sessions (split conformal, CQR), coverage measured on test sessions, overall and per volatility decile. & Coverage of \PICPEnsMin{} to \PICPEnsMax{} measured on unseen sessions. The coverage gap across volatility deciles shrinks on four links (Section~\ref{sec:why}). \\
Model output & One value per target and step. & Mean and five non-crossing quantiles (DS2S), anchored decoder, causal context features. & 95\% interval score $-$\BSISWorstAbs\% to $-$\BSISBestAbs\%, RTT MAE better on five links, RMSE within \BSRMSEMaxAbs\% (Table~\ref{tab:boot}). \\
Seeds and statistics & One training run per model, window-level summaries. & Three seeds and their ensembles, session-level bootstrap, leave-one-session-out, Wilcoxon tests. & Every interval-score gain is significant at the session level. Window-level intervals would be \BSWidRatioLo{} to \BSWidRatioHi{} times too narrow (Section~\ref{sec:validity}). \\
References & Learned models and a lag-1 ``relative RMSE''. & Persistence, seasonal repeat, ridge, linear quantile model (LQR), Gaussian heads (GS2S), the same network with a scalar output (S2S-AE), asymmetric and volatility-scaled conformal intervals. & The extent of the gain is bounded: LQR recovers most of the interval gain, DS2S adds better point forecasts (Section~\ref{sec:why}). \\
Understanding & Accuracy tables. & Analysis of what is predictable, decomposition of the interval gain, attention and gradient analysis. & Explains where the gain comes from and where it does not (Sections~\ref{sec:predictable} and~\ref{sec:attention}). \\
Use in scheduling & Argued from forecast accuracy. & Emulated deadline-aware duplication with every setting fixed on held-out sessions. & \SchedSaveSLDL\% less duplication on Starlink DL with no more late packets, ties on two links, a loss on OneWeb DL (Section~\ref{sec:sched}). \\
Reproducibility & Partial. & All code, trained models, result files, and figure scripts in the repository. Every number in the text is generated by a script. & One command per figure and table (Appendix~\ref{app:repro}). \\
\bottomrule
\end{tabularx}
\end{table}

\paragraph{Headline results.} Against the previous model retrained on the corrected data (S2S-E, three seeds), the final model (DS2S-E) improves the 95\% RTT interval score on every link, from \BSISWorstAbs\% on Starlink UL to \BSISBestAbs\% on Starlink DL, and it has the lower interval score in each of the \BSnSessTot{} test sessions. The RTT MAE improves on five links, by up to \BSBestMAEAbs\% on 5G UL. The gain comes from the calibrated distributional output: it narrows the interval when a link is quiet, widens it when the link is volatile, and places it above the forecast, where RTT spikes go (Section~\ref{sec:why}). In an emulated deadline-aware scheduler, this saves \SchedSaveSLDL\% of the duplicated packets on Starlink DL without more late packets.

\paragraph{What did not improve.} The RTT RMSE changes by at most \BSRMSEMaxAbs\% on any link. Section~\ref{sec:predictable} shows why: the level of the next RTT is close to unpredictable from a 6\,s window. On OneWeb DL the anchored decoder ranks the rare spikes worse than the previous model, which also costs duplication in the scheduler emulation. DS2S-E under-covers the most volatile tenth of OneWeb windows, and Starlink UL shows a shift between calibration and test sessions that no method corrects. A linear quantile model calibrated in the same way comes close to the final model on interval score, so the main lesson of this report is the calibrated distributional output, not the network.
